\clearpage\SourceLeaf{358}

\begin{center}
\large\bfseries [DIE KUGEL.]
\end{center}

\begin{center}
[1.]
\end{center}

\begin{center}
\bfseries ORTHOGRAPHISCHE PROJECTION DER KUGEL.
\end{center}

\begin{center}
Veränderte Erscheinung eines unendlich kleinen Theils der Kugelfläche,\\
wenn die Projectionsebene verändert wird.
\end{center}

Man wählt zur Axe der $x$ die Durchschnittslinie der beiden Projectionsebenen oder vielmehr als Axe der $y$ in der bleibenden Projectionsebene eine Parallele mit der geraden Linie, welche die beiden Projectionsplätze verbindet.

Man fragt, wie dasjenige Stück der Kugelfläche, welches in der ersten Projection wie ein Kreis erschien, in der zweiten erscheinen wird.

Coordinaten des Centrums dieser Fläche (Radius $= 1$):

\begin{center}
\begin{tabular}{ccccc}
in & der & ersten & Projection & $x,y,z$\\
,, & ,, & zweiten & ,, & $x',y',z'$.
\end{tabular}
\end{center}

Der Drehungswinkel $= \lambda$.

Es ist
\[
x' = x
\]
\[
\BareRoot{(1-xx)} = \rho
\]
\[
\BareRoot{(1-x'x')} = \rho
\]
\[
\frac{y}{\rho} = \sin u \qquad \frac{y'}{\rho} = \sin(u+\lambda)
\]
\[
dy' = \frac{z'}{z}\cdot dy - \frac{x\sin\lambda}{z}\cdot dx.
\]
